\documentclass[12pt,a4paper]{article}

% Margin configurations
\usepackage[left=2cm,top=1.25cm,right=1.25cm,bottom=1.25cm]{geometry}

% Packages for code formatting and images
\usepackage{listings}
\usepackage{xcolor}
\usepackage{graphicx}
\usepackage{hyperref}
\usepackage{float}

% Code listing style configuration
\lstset{
    basicstyle=\ttfamily\scriptsize,
    keywordstyle=\color{blue},
    stringstyle=\color{red},
    commentstyle=\color{green!60!black},
    numbers=left,
    numberstyle=\tiny\color{gray},
    stepnumber=1,
    numbersep=8pt,
    backgroundcolor=\color{gray!5},
    showspaces=false,
    showstringspaces=false,
    showtabs=false,
    frame=single,
    tabsize=4,
    captionpos=b,
    breaklines=true,
    breakatwhitespace=false,
    escapeinside={\%*}{*)}
}

\begin{document}

\begin{center}
    {\Huge \textbf{Programming in Python- Project}}\\[1cm]
    {\LARGE \textbf{Live Dynamic Weather Application}}\\[1.5cm]
\end{center}

\section*{Aim}
To design and develop a full-stack live weather application using Python (FastAPI) and a modern frontend to fetch and dynamically display real-time weather conditions alongside a 10-day predictive forecast.

\section*{Objectives}
\begin{itemize}
    \item To utilize Python's \texttt{FastAPI} and asynchronous programming concepts to create a high-performance backend server.
    \item To integrate with the Open-Meteo external REST APIs for gathering live geocoding data and weather metrics.
    \item To build an interactive, responsive frontend using HTML, CSS, and Vanilla JavaScript.
    \item To implement dynamic, pure CSS particle animations (e.g., rain, snow, drifting clouds, lightning) that respond visually to the current weather condition fetched by the Python backend.
\end{itemize}

\section*{Theory}
The project relies on a client-server architecture:
\begin{itemize}
    \item \textbf{Backend (Python):} FastAPI acts as the backend server framework. It exposes routing endpoints, including an autocomplete \texttt{/api/search} API and an \texttt{/api/weather} API. It utilizes the asynchronous \texttt{httpx} library to make concurrent external HTTP requests to Open-Meteo without blocking the main thread.
    \item \textbf{Frontend:} A static single-page application structure. The DOM is manipulated via Vanilla JavaScript. User input triggers asynchronous \texttt{fetch()} calls to the Python backend. The presentation layer utilizes UI principles like glassmorphism and keyframe animations to create immersive environment systems for weather visualization.
\end{itemize}

\section*{Code}

\subsection*{1. Backend Configuration (main.py)}
\begin{lstlisting}[language=Python, caption=Backend Logic (main.py)]
from fastapi import FastAPI, HTTPException, Query
from fastapi.staticfiles import StaticFiles
from fastapi.responses import FileResponse
import httpx
import logging

app = FastAPI(title="Live Weather App")
logger = logging.getLogger(__name__)

def get_weather_condition(code: int) -> str:
    if code in [0, 1]: return "clear"
    elif code in [2, 3, 45, 48]: return "clouds"
    elif code in [51, 53, 55, 56, 57, 61, 63, 65, 66, 67, 80, 81, 82]: return "rain"
    elif code in [71, 73, 75, 77, 85, 86]: return "snow"
    elif code in [95, 96, 99]: return "thunderstorm"
    return "clear"

@app.get("/api/search")
async def search_city(q: str = Query(..., min_length=1)):
    async with httpx.AsyncClient() as client:
        geocode_url = f"https://geocoding-api.open-meteo.com/v1/search?name={q}&count=5"
        resp = await client.get(geocode_url)
        data = resp.json()
        results = []
        if "results" in data:
            for r in data["results"]:
                results.append({
                    "id": r.get("id"),
                    "name": r.get("name"),
                    "country": r.get("country"),
                    "admin1": r.get("admin1"),
                    "lat": r["latitude"],
                    "lon": r["longitude"]
                })
        return {"results": results}

@app.get("/api/weather")
async def get_weather(lat: float = None, lon: float = None, city: str = None):
    if lat is None or lon is None:
        if not city:
            raise HTTPException(status_code=400, detail="Must provide either lat/lon or city name")
        # Fallback to geocode if only city name provided
        async with httpx.AsyncClient() as client:
            geocode_url = f"https://geocoding-api.open-meteo.com/v1/search?name={city}&count=1"
            geo_response = await client.get(geocode_url)
            geo_data = geo_response.json()
            if not geo_data.get("results"):
                raise HTTPException(status_code=404, detail="City not found")
            location = geo_data["results"][0]
            lat = location["latitude"]
            lon = location["longitude"]
            city_name = location["name"]
            country = location.get("country", "")
    else:
        city_name = city or "Selected Location"
        country = ""

    async with httpx.AsyncClient() as client:
        weather_url = (
            f"https://api.open-meteo.com/v1/forecast?latitude={lat}&longitude={lon}"
            "&current=temperature_2m,relative_humidity_2m,is_day,precipitation,weather_code,wind_speed_10m"
            "&daily=weather_code,temperature_2m_max,temperature_2m_min"
            "&timezone=auto&forecast_days=10"
        )
        weather_response = await client.get(weather_url)
        weather_data = weather_response.json()
        
        if "current" not in weather_data or "daily" not in weather_data:
            raise HTTPException(status_code=500, detail="Weather data unavailable")
            
        current = weather_data["current"]
        daily = weather_data["daily"]
        
        forecast = []
        for i in range(len(daily["time"])):
            forecast.append({
                "date": daily["time"][i],
                "max_temp": daily["temperature_2m_max"][i],
                "min_temp": daily["temperature_2m_min"][i],
                "condition": get_weather_condition(daily["weather_code"][i]),
                "weather_code": daily["weather_code"][i]
            })

        return {
            "city": city_name,
            "country": country,
            "temperature": current.get("temperature_2m"),
            "humidity": current.get("relative_humidity_2m"),
            "wind_speed": current.get("wind_speed_10m"),
            "precipitation": current.get("precipitation"),
            "is_day": current.get("is_day") == 1,
            "weather_code": current.get("weather_code"),
            "condition": get_weather_condition(current.get("weather_code", 0)),
            "forecast": forecast
        }

@app.get("/")
def serve_index():
    return FileResponse("static/index.html")

app.mount("/", StaticFiles(directory="static"), name="static")

if __name__ == "__main__":
    import uvicorn
    uvicorn.run("main:app", host="127.0.0.1", port=8000, reload=True)
\end{lstlisting}

\subsection*{2. Frontend HTML Structure (index.html)}
\begin{lstlisting}[language=HTML, caption=Frontend Structure (static/index.html)]
<!DOCTYPE html>
<html lang="en">
<head>
    <meta charset="UTF-8">
    <meta name="viewport" content="width=device-width, initial-scale=1.0">
    <title>Live Weather Application</title>
    <link rel="stylesheet" href="style.css">
    <link href="https://fonts.googleapis.com/css2?family=Inter:wght@300;400;500;600;700&display=swap" rel="stylesheet">
</head>
<body>
    <div id="app-background" class="bg-landing">
        <div id="bg-layer-stars" class="bg-layer hidden"></div>
        <div id="bg-layer-sun" class="bg-layer hidden"></div>
        <div id="bg-layer-clouds" class="bg-layer hidden"></div>
        <div id="bg-layer-rain" class="bg-layer hidden"></div>
        <div id="bg-layer-snow" class="bg-layer hidden"></div>
        <div id="bg-layer-lightning" class="bg-layer hidden"></div>
    </div>

    <div id="app-wrapper" class="state-landing">
        <div class="search-section">
            <div id="landing-hero" class="landing-hero">
                <h1>Global Weather Insights.</h1>
                <p>Live conditions and 10-day forecasts.</p>
            </div>
            
            <div class="search-container">
                <div class="search-box">
                    <input type="text" id="city-input" placeholder="Search for a city..." autocomplete="off">
                    <button id="search-btn" aria-label="Search">
                        <svg viewBox="0 0 24 24" width="20" height="20" stroke="currentColor" stroke-width="2" fill="none" stroke-linecap="round" stroke-linejoin="round" class="css-i6dzq1"><circle cx="11" cy="11" r="8"></circle><line x1="21" y1="21" x2="16.65" y2="16.65"></line></svg>
                    </button>
                    <div id="search-spinner" class="spinner hidden"></div>
                </div>
                <div id="suggestions" class="suggestions-dropdown hidden"></div>
            </div>
        </div>
        
        <div id="content-section" class="content-section hidden">
            <div id="loading" class="status-msg hidden">Gathering atmospheric data...</div>
            <div id="error-msg" class="status-msg hidden"></div>
            
            <div id="weather-scroll-area" class="scroll-area hidden">
                <div class="cards-layout">
                    <!-- Current Weather Card -->
                    <div id="weather-card" class="weather-card panel">
                        <div class="location-header">
                            <h1 id="city-name">--</h1>
                            <p id="country-name">--</p>
                        </div>
                        <div class="main-weather">
                            <div class="weather-icon-wrapper">
                                <div id="weather-icon-slot"></div>
                            </div>
                            <div class="temp-wrapper">
                                <span id="temperature">--</span><span class="celsius">&deg;</span>
                            </div>
                            <p id="condition-text" class="condition-description">--</p>
                        </div>
                        <div class="weather-details">
                            <div class="detail-item">
                                <span class="label">Humidity</span><span class="value" id="humidity">--%</span>
                            </div>
                            <div class="detail-item">
                                <span class="label">Wind</span><span class="value" id="wind">-- km/h</span>
                            </div>
                            <div class="detail-item">
                                <span class="label">Precip</span><span class="value" id="precip">-- mm</span>
                            </div>
                        </div>
                    </div>
    
                    <!-- 10 Day Forecast Card -->
                    <div id="forecast-card" class="forecast-card panel">
                        <h3>10-Day Forecast</h3>
                        <div id="forecast-list" class="forecast-list"></div>
                    </div>
                </div>
            </div>
        </div>
    </div>
    <script src="app.js"></script>
</body>
</html>
\end{lstlisting}

\subsection*{3. Frontend JavaScript Logic (app.js)}
\begin{lstlisting}[caption=Frontend Logic (static/app.js)]
document.addEventListener("DOMContentLoaded", () => {
    const searchBtn = document.getElementById("search-btn");
    const cityInput = document.getElementById("city-input");
    const suggestionsBox = document.getElementById("suggestions");
    const searchSpinner = document.getElementById("search-spinner");
    
    const appWrapper = document.getElementById("app-wrapper");
    const contentSection = document.getElementById("content-section");
    const weatherScrollArea = document.getElementById("weather-scroll-area");
    const loading = document.getElementById("loading");
    const errorMsg = document.getElementById("error-msg");
    const appBackground = document.getElementById("app-background");
    
    let debounceTimer;

    cityInput.addEventListener("input", () => {
        clearTimeout(debounceTimer);
        const query = cityInput.value.trim();
        if (query.length < 2) {
            suggestionsBox.classList.add("hidden");
            searchSpinner.classList.add("hidden");
            return;
        }
        searchSpinner.classList.remove("hidden");
        debounceTimer = setTimeout(() => { fetchSuggestions(query); }, 400);
    });

    cityInput.addEventListener("keypress", (e) => {
        if (e.key === "Enter" && cityInput.value.trim() !== "") {
            suggestionsBox.classList.add("hidden");
            fetchWeather(null, null, cityInput.value.trim());
        }
    });

    searchBtn.addEventListener("click", () => {
        if (cityInput.value.trim() !== "") {
            suggestionsBox.classList.add("hidden");
            fetchWeather(null, null, cityInput.value.trim());
        }
    });

    document.addEventListener("click", (e) => {
        if (!e.target.closest('.search-container')) suggestionsBox.classList.add("hidden");
    });

    async function fetchSuggestions(query) {
        try {
            const res = await fetch(`/api/search?q=${encodeURIComponent(query)}`);
            const data = await res.json();
            searchSpinner.classList.add("hidden");
            if (data.results && data.results.length > 0) renderSuggestions(data.results);
            else suggestionsBox.classList.add("hidden");
        } catch (err) { searchSpinner.classList.add("hidden"); }
    }

    function renderSuggestions(results) {
        suggestionsBox.innerHTML = "";
        results.forEach(res => {
            const div = document.createElement("div");
            div.className = "suggestion-item";
            const admin = res.admin1 ? `${res.admin1}, ` : "";
            const country = res.country || "";
            div.innerHTML = `<span class="sugg-city">${res.name}</span><span class="sugg-region">${admin}${country}</span>`;
            div.addEventListener("click", () => {
                cityInput.value = res.name;
                suggestionsBox.classList.add("hidden");
                fetchWeather(res.lat, res.lon, res.name);
            });
            suggestionsBox.appendChild(div);
        });
        suggestionsBox.classList.remove("hidden");
    }

    async function fetchWeather(lat, lon, city) {
        appWrapper.className = "state-weather";
        contentSection.classList.remove("hidden");
        weatherScrollArea.classList.add("hidden");
        errorMsg.classList.add("hidden");
        loading.classList.remove("hidden");
        cityInput.blur();

        try {
            let url = `/api/weather?`;
            if (lat && lon) url += `lat=${lat}&lon=${lon}&city=${encodeURIComponent(city)}`;
            else url += `city=${encodeURIComponent(city)}`;
            
            const res = await fetch(url);
            const data = await res.json();
            if (!res.ok) throw new Error(data.detail || "Unable to fetch weather");
            updateUI(data);
        } catch (err) {
            loading.classList.add("hidden");
            errorMsg.textContent = err.message;
            errorMsg.classList.remove("hidden");
        }
    }

    function updateUI(data) {
        loading.classList.add("hidden");
        weatherScrollArea.classList.remove("hidden");

        document.getElementById("city-name").textContent = data.city;
        document.getElementById("country-name").textContent = data.country || "Weather Report";
        document.getElementById("temperature").textContent = Math.round(data.temperature);
        document.getElementById("humidity").textContent = `${data.humidity}%`;
        document.getElementById("wind").textContent = `${data.wind_speed} km/h`;
        document.getElementById("precip").textContent = `${data.precipitation || 0} mm`;

        let condition = data.condition;
        let isDay = data.is_day;
        let displayCondition = condition;
        if(condition === "clear" && !isDay) displayCondition = "Clear Night";
        if(condition === "clear" && isDay) displayCondition = "Sunny";
        
        document.getElementById("condition-text").textContent = displayCondition;
        updateThemeAndIcon(condition, isDay);
        updateForecast(data.forecast);
    }

    function updateThemeAndIcon(condition, isDay) {
        let bgClass = "bg-clear-day";
        if (condition === "clear" || condition === "clouds") bgClass = `bg-${condition}-${isDay ? 'day' : 'night'}`;
        else bgClass = `bg-${condition}`;
        appBackground.className = bgClass;
        
        const iconSlot = document.getElementById("weather-icon-slot");
        let iconHTML = "";
        switch (condition) {
            case "clear": iconHTML = isDay ? `<div class="icon-sun"></div>` : `<div class="icon-moon"></div>`; break;
            case "clouds": iconHTML = `<div class="icon-cloud"></div>`; break;
            case "rain": iconHTML = `<div class="icon-rain-cloud"><div class="raindrop"></div><div class="raindrop"></div><div class="raindrop"></div></div>`; break;
            case "snow": iconHTML = `<div class="icon-snow-cloud"><div class="snowflake"></div><div class="snowflake"></div><div class="snowflake"></div></div>`; break;
            case "thunderstorm": iconHTML = `<div class="icon-thunder-cloud"><div class="lightning"></div></div>`; break;
            default: iconHTML = isDay ? `<div class="icon-sun"></div>` : `<div class="icon-moon"></div>`;
        }
        iconSlot.innerHTML = iconHTML;
        renderBackgroundEffects(condition, isDay);
    }

    function renderBackgroundEffects(condition, isDay) {
        const layers = document.querySelectorAll(".bg-layer");
        layers.forEach(l => { l.innerHTML = ""; l.classList.add("hidden"); });

        if (condition === "clear" && !isDay) {
            const stars = document.getElementById("bg-layer-stars");
            stars.classList.remove("hidden");
            for(let i=0; i<50; i++) {
                let div = document.createElement("div"); div.className = "star";
                div.style.left = `${Math.random()*100}vw`; div.style.top = `${Math.random()*60}vh`; 
                div.style.width = div.style.height = `${Math.random()*3 + 1}px`;
                div.style.animationDuration = `${Math.random()*2 + 1}s`;
                stars.appendChild(div);
            }
        }
        else if (condition === "rain" || condition === "thunderstorm") {
            const rain = document.getElementById("bg-layer-rain");
            rain.classList.remove("hidden");
            for(let i=0; i<60; i++) {
                let div = document.createElement("div"); div.className = "bg-drop";
                div.style.left = `${Math.random()*100}vw`; div.style.height = `${Math.random()*20 + 20}px`;
                div.style.animationDuration = `${Math.random()*0.5 + 0.5}s`;
                rain.appendChild(div);
            }
        }
    }

    function updateForecast(forecast) {
        const list = document.getElementById("forecast-list");
        list.innerHTML = "";
        forecast.forEach((day, index) => {
            const d = new Date(day.date);
            const dateStr = index === 0 ? "Today" : index === 1 ? "Tomorrow" : d.toLocaleDateString("en-US", {weekday: 'short', month: 'short', day: 'numeric'});
            let icon = day.condition === "clear" ? "[Sun]" : day.condition === "clouds" ? "[Cloud]" : day.condition === "rain" ? "[Rain]" : day.condition === "snow" ? "[Snow]" : "[Storm]";
            const div = document.createElement("div");
            div.className = "forecast-item";
            div.innerHTML = `<div class="fc-date">${dateStr}</div><div class="fc-icon">${icon}</div>
                             <div class="fc-temps"><span class="fc-max">${Math.round(day.max_temp)}&deg;</span><span class="fc-min">${Math.round(day.min_temp)}&deg;</span></div>`;
            list.appendChild(div);
        });
    }
});
\end{lstlisting}

\section*{Output}
\begin{figure}[H]
    \centering
    \includegraphics[width=0.55\textwidth]{landing_page.png}
    \caption{Landing Page}
\end{figure}

\begin{figure}[H]
    \centering
    \includegraphics[width=0.55\textwidth]{smart_search.png}
    \caption{Smart autocomplete suggestions working perfectly across location searches.}
\end{figure} 

\begin{figure}[H]
    \centering
    \includegraphics[width=0.55\textwidth]{real_time_weather.png}
    \caption{The application displaying real-time weather parameters and animations.}
\end{figure} 

\clearpage

\begin{figure}[H]
    \centering
    \includegraphics[width=0.65\textwidth]{10_day_forecast.png}
    \caption{10-day predictive forecast card dynamically fetched from the backend API.}
\end{figure}


\begin{itemize}
    \item \textbf{Test Case 1:} Searching for "Tokyo" successfully fetched atmospheric data and populated the 10-day forecast card.
    \item \textbf{Test Case 2:} Validating autocomplete suggestions by calling the Python backend while typing city prefixes.
\end{itemize}

\section*{Learning Outcomes}
\begin{itemize}
    \item Gained practical experience in building RESTful APIs using Python's modern FastAPI framework.
    \item Understood the mechanics of asynchronous HTTP requests in Python using the \texttt{httpx} library.
    \item Built knowledge on connecting backend systems to third-party open APIs (Open-Meteo).
    \item Strengthened web-development concepts including asynchronous JavaScript (\texttt{async/await}), cross-origin resource handling, and advanced CSS styling (keyframe animations, glassmorphism layouts).
\end{itemize}

\end{document}
